% ---- pg-264 （书页 252） ----

阵 $\boldsymbol{\Sigma}$．与一元方法不同的是，该模型并不要求四次测量具有相同的方差．由 4 个样本均值 $(\bar{x}_1,\bar{x}_2,\bar{x}_3,\bar{x}_4)$ 构成的轮廓概要地描述了特征的生长过程，即钙随时间损失的过程．问题是，这种生长模式是否可用时间的多项式来表达？

\begin{mvtable}
\phantomsection\label{tab:6.5}\mvtabcap{表 6.5\quad 被测对象主尺骨中的钙含量}
\begin{tabular}{crrrr}
\toprule
被测个体 & 起初 & 1 年后 & 2 年后 & 3 年后\\
\midrule
1 & 87.3 & 86.9 & 86.7 & 75.5\\
2 & 59.0 & 60.2 & 60.0 & 53.6\\
3 & 76.7 & 76.5 & 75.7 & 69.5\\
4 & 70.6 & 76.1 & 72.1 & 65.3\\
5 & 54.9 & 55.1 & 57.2 & 49.0\\
6 & 78.2 & 75.3 & 69.1 & 67.6\\
7 & 73.7 & 70.8 & 71.8 & 74.6\\
8 & 61.8 & 68.7 & 68.2 & 57.4\\
9 & 85.3 & 84.4 & 79.2 & 67.0\\
10 & 82.3 & 86.9 & 79.4 & 77.4\\
11 & 68.6 & 65.4 & 72.3 & 60.8\\
12 & 67.8 & 69.2 & 66.3 & 57.9\\
13 & 66.2 & 67.0 & 67.0 & 56.2\\
14 & 81.0 & 82.3 & 86.8 & 73.9\\
15 & 72.3 & 74.6 & 75.3 & 65.1\\
\midrule
均值 & 72.38 & 73.29 & 72.47 & 64.79\\
\bottomrule
\end{tabular}
\end{mvtable}

资料来源：数据经 Everett Smith 许可使用．

在时刻 $t_1,t_2,\cdots,t_p$，对所有个体进行 $p$ 次测量后，波托夫-罗伊的二次生长模型就变为
\[
E[\mathbf{X}]=E\begin{bmatrix}X_1\\X_2\\\vdots\\X_p\end{bmatrix}
=\begin{bmatrix}\beta_0+\beta_1t_1+\beta_2t_1^2\\ \beta_0+\beta_1t_2+\beta_2t_2^2\\ \vdots\\ \beta_0+\beta_1t_p+\beta_2t_p^2\end{bmatrix}
\]
其中第 $i$ 个均值为 $t_i$ 的二次函数．

通常需要对各组进行比较．\hyperref[tab:6.6]{表 6.6} 给出第二组（姑称之为处理组）妇女（即接受一项饮食调节和定期锻炼计划的老年妇女）主尺骨中的钙含量测量值．

\begin{mvtable}
\phantomsection\label{tab:6.6}\mvtabcap{表 6.6\quad 主尺骨中的钙含量（处理组）}
\begin{tabular}{crrrr}
\toprule
被测个体 & 起初 & 1 年后 & 2 年后 & 3 年后\\
\midrule
1 & 83.8 & 85.5 & 86.2 & 81.2\\
2 & 65.3 & 66.9 & 67.0 & 60.6\\
3 & 81.2 & 79.5 & 84.5 & 75.2\\
4 & 75.4 & 76.7 & 74.3 & 66.7\\
5 & 55.3 & 58.3 & 59.1 & 54.2\\
6 & 70.3 & 72.3 & 70.6 & 68.6\\
7 & 76.5 & 79.9 & 80.4 & 71.6\\
8 & 66.0 & 70.9 & 70.3 & 64.1\\
\bottomrule
\end{tabular}
\end{mvtable}
